\documentclass[11pt]{exam}
\usepackage[empty]{fullpage}
\usepackage{amsmath}
\usepackage{tikz}
\usepackage{pgfplots}
\usepackage{hyperref}
\usepackage{soul}

\newcommand{\ds}{\displaystyle}
\newcommand{\on}{\operatorname}
\newcommand{\Var}{\operatorname{Var}}

\begin{document}
\subsection*{Midterm 3 Review \hfill Math 222}

\textit{These are problems similar to the ones that might be on midterm 3. Be sure to also review class workshops and problems from the textbook.}

\begin{questions}
\question A 2003 study (North, Shilcock, \& Hargreaves) looked at whether background music in a restaurant affected the amount of money that diners spent on their meals. The researchers asked a restaurant to alternate classical music, popular music, and silence on successive nights over 18 days. The following summary statistics were reported:

\begin{center}
\begin{tabular}{|l|c|c|c|}
\hline
~ & Classical Music & Pop Music & No Music \\ \hline
Mean & \pounds24.13 & £21.91 & £21.70 \\ 
SD & £2.243 & £2.627 & £3.332 \\ 
Sample size & 120 & 142 & 131 \\ \hline
\end{tabular}
\end{center}
\begin{parts}
\part There were a total of 393 customers at the restaurant.  What was the overall average amount of money spent by these customers? 

\part Use a calculator to find the SSG and SSE for this data, and then make an ANOVA table for this data including the F-statistic.   

\part Use the R command \verb|1 - pf(x, df1, df2)| to write down an expression for the p-value in this situation. 

\part What additional information would you need before you can safely draw conclusions from this p-value?

\part Why might the independence assumption not be satisfied in this example? 

\end{parts}

\question OpenIntro (\url{https://people.hsc.edu/faculty-staff/blins/books/OpenIntroStats4e.pdf#eoce.7.39}) is a good exercise.

\end{questions}


\end{document}
